\section{A Gaussian-query proposal}\label{sec:proposal}

We now construct a spin using the Gaussian matrix and auxiliary
randomness. Each matrix query is orthogonal to the earlier queries.
This exposes fresh Gaussian information while leaving a compressed GOE
block unrevealed. The scalar identity of Section~\ref{sec:scalar} will
then control both the relative entropy of the resulting law and its
likelihood outside a small exceptional event.

\subsection{Construction and conditional Gaussian laws}

Successive Gaussian revelations along orthogonalized directions occur
in the iterative TAP construction~\cite[Sections~2--3]{Bolthausen2014}.
Incremental approximate message passing uses related Parisi-guided
normalized increments with asymptotic orthogonality
\cite[Sections~2--3]{Montanari2021}. Here we give an exact
conditional-law argument for our finite mesh, including the auxiliary
noises that complete each increment.

Fix $0<l\le1/100$ and define
\[
 h=n^{-1/16},\qquad K=\lceil l\log n/h\rceil,\qquad
 t_k=kh,\qquad V=t_K.
\]
All assertions below are for sufficiently large $n$, so $K<n$ and
$V>v_*$. The disorder $G$ has its ordinary Gaussian law. We construct
a vector process with coordinates following the scalar SDE, using
\[
 dZ=\rho(v)\,dB+\sqrt{1-\rho(v)^2}\,dD.
 \tag{4}\label{eq:4}
\]
Here $D$ is an independent auxiliary Brownian motion. The stream $B$
is constructed from matrix queries and auxiliary completions so that,
after averaging over the disorder, $(B,D)$ has independent standard
Brownian coordinates. The weight $\rho$ equals one up to $v_*$ and
is nonincreasing afterward. The continuation therefore reduces the
share of the query-generated stream while keeping unit variance in
$Z$. Put $M(v)_i=m(v,Y_{v,i})$.

On $[0,h]$, generate $B$ entirely from independent auxiliary Brownian
noise. For each subsequent interval $[t_k,t_{k+1}]$, $1\le k<K$,
the previous scalar increment determines
\[
 s_k^2=\int_{t_{k-1}}^{t_k}r(v)\,dv,
 \qquad T_k=\frac{M(t_k)-M(t_{k-1})}{s_k}.
\]
Project $T_k$ off the preceding query directions $u_1,\ldots,u_{k-1}$,
and normalize the residual to obtain $|u_k|_n=1$. If the residual is
zero, choose a new direction measurably from the past: project the
coordinate basis vectors in their fixed order, take the first nonzero
projection, and normalize it. Since $k<n$, such a direction exists.
Let $\Pi_{\le k}$ be the orthogonal projection onto the span through
$u_k$, and write $\Pi_{>k}=I-\Pi_{\le k}$.

Query the column $Wu_k$. For an independent standard Gaussian vector
$z_k$, set the endpoint of the next Brownian increment to be
\[
 g_k=h^{-1/2}\bigl(B(t_{k+1})-B(t_k)\bigr)
     =\Pi_{>k}Wu_k+\Pi_{\le k}z_k.
\]
Fill the interval with linear interpolation between its endpoints plus
an independent centered Brownian bridge. Every new bridge and completion
noise is fresh. Deterministic-coefficient integrals against this filled
path may equivalently be obtained by conditioning Brownian motion on
its endpoint. At time $V$, draw the spins independently given the path,
with means $M_i(V)$. Call the marginal law of $(G,x)$ obtained this way
$Q_l$.

To verify the stated Brownian law, condition on the previously revealed
columns and the auxiliary inputs already used. Gaussian orthogonal
invariance leaves the unqueried block GOE on the orthogonal complement,
with the original variance scale. The next direction is measurable with
respect to this past. In an orthonormal basis beginning with
$u_k/\sqrt n$, the new diagonal entry, the remaining entries of that
column, and the next compressed block are independent. Therefore
$\Pi_{>k}Wu_k$ is a fresh Gaussian with covariance $\Pi_{>k}$, and
\[
 d_k=\langle u_k,Wu_k\rangle_n\sim N(0,2/n)
\]
conditionally before the column is revealed. The independent completion
has covariance $\Pi_{\le k}$, so $g_k$ is a fresh standard Gaussian
vector. Induction, followed by the independent bridge filling, proves
the assertion about $(B,D)$ in \eqref{eq:4}.

In particular, after forgetting the matrix, the coordinate paths and
their final spins are independent and identically distributed. The
conditional mean of $x$ given the Brownian streams through time $v$ is
$M(v)$. Its unconditional law is uniform on the cube by the symmetry of
the scalar process. The spin marginal of $P$ is also uniform, by
coordinate-sign symmetry of the disorder. These facts concern marginals;
the proposal spin remains correlated with its matrix.

\subsection{Entropy and likelihood estimates}

The following proposition supplies two distinct comparisons. Entropy
will yield coverage after smoothing. The likelihood estimate will keep
intermediate proposals out of the small boundary of a putative low-flow
cut.

\begin{proposition}[Proposal comparison]\label{prop:proposal}
For every fixed $0<l\le1/100$ and all sufficiently large $n$,
\[
 {\rm KL}(Q_l\Vert P)\le2n^{1-l/2}.
 \tag{5}\label{eq:5}
\]
There is an event in the proposal experiment of probability at least
$1-\exp(-n^{1-10l})$ such that the output subprobability restricted to
this event is dominated, on the space of $(G,x)$, by
\[
 \exp(n^{1-l/2})P.
 \tag{6}\label{eq:6}
\]
The empirical Gram matrix of the $T_k$ is within $n^{-l}$ of the
identity in operator norm, except with probability at most
$\exp(-n^{1-10l})$. The dimension required for these estimates may
depend on the fixed $l$.
\end{proposition}

\begin{proof}
Put $\eta=n^{-l}$. We first evaluate the proposal energy from empirical
inner products, then compute the information cost of its matrix
channel. The pressure bound and \eqref{eq:3} combine these estimates.

\paragraph{Empirical inner products.}
Form the empirical matrix of all normalized inner products among
$T_k,g_k$ ($1\le k<K$) and $x$. It is an average of $n$ iid row
outer products, by the coordinate law just proved. Every unit linear
form of a row has sub-Gaussian norm at most $C_\beta e^V$.
For the $T_k$ part, use their disjoint martingale increments,
$|a|\le1$, and
\[
 s_k^2\ge c_\beta h e^{-2V},
\]
which follows from \eqref{eq:1}. An exponential martingale estimate
then bounds the linear form by its deterministic quadratic-variation
bound. The $g_k$ part is Gaussian and the spin term is bounded; their
correlations with the first part do not affect the resulting
sub-Gaussian bound.

For completeness, squared linear forms have moments bounded by
$k!(C_\beta e^{2V})^k$. Expanding the centered exponential moment
and applying exponential Markov gives, at deviation scale $\eta/2$,
a quadratic-form tail bounded by
\[
 2\exp(-c_\beta n\eta^2e^{-4V}).
\]
A $1/4$-net of the unit sphere in dimension $O(K)$ has
$\exp(O(K))$ points, by ball packing. The quadratic-form norm
inequality and a union bound therefore put the entire empirical
matrix within $\eta$ of its expectation, with exceptional probability
at most $\exp(-n^{1-8l})$.

The population identities are
\[
 \begin{gathered}
 \mathbb E[T_jT_k]_{\rm row}=\mathbf1_{j=k},\qquad
 \mathbb E[T_jg_k]_{\rm row}=0\quad(j\le k),\\
 \mathbb E[T_kx]_{\rm row}=s_k,\qquad
 \mathbb E[g_kx]_{\rm row}=b_k
   :=h^{-1/2}\int_{t_k}^{t_{k+1}}p(v)\rho(v)\,dv.
 \end{gathered}
\]
They follow from the martingale bracket, the conditional mean of the
final spin, and \eqref{eq:4}. The vectors $s=(s_k)$ and $b=(b_k)$
each have Euclidean norm at most one, by projection of the unit-variance
spin onto the corresponding orthonormal scalar variables.

We need matrix-norm estimates that do not lose a factor $K$ under
orthogonalization. Taking the triangular part of an $m\times m$
matrix has operator-norm cost at most $C\log(2m)$. To see this, pad
to a power of two and partition the strict triangle by dyadic scales.
At each scale the rectangles have disjoint row and column sets, so
retaining that scale is contractive; summing over scales proves the
bound.

Let $A$ be a Gram matrix with $\|A-I\|\le\eta<1/2$, and let
$R^{\mathsf T}R=A$ be its upper Cholesky factor. Along
$A(t)=I+t(A-I)$, differentiation shows that
$\dot R R^{-1}$ is the upper triangular part, with half its diagonal,
of $R^{-\mathsf T}\dot A R^{-1}$. The norms of $R$ and $R^{-1}$
are uniformly bounded. The preceding triangular estimate and
integration give
$\|R-I\|\le C\eta\log(2m)$. Differentiability follows directly
from the triangular Cholesky recursion.

Apply this with the Gram matrix of the $T_k$. The query matrix with
columns $u_k$ is the matrix with columns $T_k$, multiplied by $R^{-1}$.
Thus, writing $c_k=\langle u_k,x\rangle_n$, the empirical event gives
\[
 \begin{aligned}
 \|c-s\|&\le C\eta\log(2K),\\
 \left\|\bigl(\mathbf1_{j\le k}\langle u_j,g_k\rangle_n\bigr)_{jk}
       \right\|&\le C\eta\log(2K),\\
 \| (\langle g_k,x\rangle_n)_k-b\|&\le\eta.
 \end{aligned}
\]
For the middle estimate, left multiplication of the population cross
matrix by the lower triangular $R^{-\mathsf T}$ preserves its zero
upper triangle, including the diagonal. Also $\|c\|\le1$, because
the query directions are orthonormal in the normalized inner product.

\paragraph{Energy of the final spin.}
Let $H_0=\Pi_{>K-1}$. Revealing the successive columns decomposes the
quadratic form exactly as
\[
 \frac{\widehat H_G(x)}n
 =\sum_kc_k\langle x,\Pi_{>k}g_k\rangle_n
  +\frac12\sum_kc_k^2d_k
  +\frac12\langle x,H_0WH_0x\rangle_n.
\]
The last two terms are $O(\eta)$ outside an event of probability at
most $CK\exp(-cn\eta^2)$. The diagonal estimate follows from the
conditional Gaussian law and $\sum c_k^2\le1$. For the residual
term, the spin has not used the unrevealed block, so its quadratic
form has the required conditional Gaussian tail.

Removing the projections in the first term costs at most
$C\eta\log(2K)$ by the triangular matrix bound above. Its remaining
value is $s\cdot b$ up to the same error. The bound on
$(\log r)'$ implies that $s_k/\sqrt h$ differs from
$\sqrt{r(v)}$ on $[t_k,t_{k+1}]$ by a relative $O(h)$.
The missing initial interval contributes $O(h)$, and $e(V)$ is
bounded. Hence $s\cdot b=e(V)+O_\beta(h)$, and altogether
\[
 \widehat H_G(x)/n=e(V)+O_\beta(h+\eta\log(2K)).
 \tag{7}\label{eq:7}
\]
This is the energy contribution to the comparison. We next isolate the
entropy paid for the matrix information alone.

\paragraph{The auxiliary reference law.}
Start on an enlarged path space with an independent uniform spin $x$,
the ordinary Gaussian disorder, and all the unconditioned proposal
inputs. Call this the unlinked experiment. To obtain the actual
proposal, multiply its density by
\[
 \prod_i(1+x_i m(V,Y_{V,i})).
\]
For fixed $x$, the scalar martingale shows that this is a stochastic
exponential with coordinate integrands
\[
 \xi_i(v)=\frac{x_i a(v,Y_{v,i})}{1+x_i m(v,Y_{v,i})},
 \qquad |\xi_i(v)|\le2,
\]
against the streams $Z_i$ in \eqref{eq:4}. Call the resulting law the
linked experiment. Summing the linking density over the uniform spin
shows that the path and disorder marginal is unchanged.

For a continuous local martingale $N$ starting at zero, write
$\mathcal E(N)_V=\exp(N_V-\langle N\rangle_V/2)$.
The linking density factors as $E_BE_D$, where
\[
 \begin{aligned}
 E_B&=\mathcal E\left(\sum_i\int_0^{\cdot}
                  \xi_i\rho\,dB_i\right)_V,\\
 E_D&=\mathcal E\left(\sum_i\int_0^{\cdot}
                  \xi_i\sqrt{1-\rho^2}\,dD_i\right)_V.
 \end{aligned}
\]
This product identity uses the zero quadratic covariation of $B$ and
$D$. Define the auxiliary reference to have density $E_D$ relative
to the unlinked law. Its $(G,x)$ marginal is still the independent product
of Gaussian disorder and uniform spin. The conditional mean-one
assertion behind this fact needs a filtration check, because the
matrix-query endpoints depend on the past of $D$.

Let $\mathcal A$ contain all primitive auxiliary inputs other than
$D$, and under the unlinked law set
\[
 \mathcal F_v=\sigma(x,B_{[0,v]},D_{[0,v]}),\qquad
 \mathcal G_v=\sigma(x,G,\mathcal A,D_{[0,v]}).
\]
Interval by interval, every query direction is determined by past
$D$ and the information in $\mathcal G_0$. Each new endpoint and
its entire bridge are then known before using the current $D$
increment. Deterministic-coefficient bridge integrals can be formed as
linear Gaussian integrals, with the endpoint term fixed by that past.
Thus $\mathcal F_v\subseteq\mathcal G_v$.
One can also see the inclusion pathwise: $\rho=1$ before $v_*$, and
\eqref{eq:1} gives $-\rho\le\rho'\le0$ afterward, so $\rho$ has
bounded variation. The identity
$\int\rho\,dB=\rho B-\int B\,d\rho$, applied on successive
intervals, uses no future $D$.

The pair $(B,D)$ is Brownian in $\mathcal F$, whereas $D$ is
Brownian in $\mathcal G$ since its time-zero information is
independent of $D$. The bounded continuous $D$ integrand has the
same stochastic integral in these two filtrations: its identical
left-endpoint approximations converge in $L^2$ by the It\^o isometry.
The bounded-integrand exponential-martingale theorem in $\mathcal G$
now gives conditional mean one given $(x,G,\mathcal A)$.
This proves the claimed marginal for the auxiliary reference.
All subsequent Brownian changes of measure use $\mathcal F$;
no Brownian property for $B$ in the enlarged filtration is required.

\paragraph{Information cost and likelihood.}
The density of the linked experiment relative to the auxiliary
reference is $E_B$. Its logarithm is
\[
 L=\log E_B=\sum_i\left(\int_0^V\xi_i\rho\,dB_i
                  -\frac12\int_0^V\xi_i^2\rho^2\,dv\right).
\]
Under the linked law, compensate the $B_i$ drift by $\rho\xi_i$.
The resulting centered integral has bracket at most $4nV$ and hence
has exponential-martingale tails. The remaining half-bracket is a sum
of independent coordinate variables bounded by $2V$ each. Coordinate
independence holds after linking because the linking density factors
across the iid coordinate streams and uniform spins. Moreover,
conditional on the streams through time $v$, the linked spin has
probabilities $(1\pm m(v,Y_{v,i}))/2$. It follows that
\[
 \mathbb E L=nj(V),\qquad L=nj(V)+O(n\eta)
\]
except with probability at most $\exp(-n^{1-4l})$ for large $n$.
The centered integral is controlled by exponential Markov; the
half-bracket by bounded-variable concentration. The Brownian drift
change itself follows either from Girsanov's theorem for bounded
integrands or from bounded predictable step approximation and the
mean-one exponential identities.

Change the auxiliary reference a second time, multiplying only by the
density of $P$ relative to independent disorder and uniform spin.
The new reference has $(G,x)$ marginal $P$, and the log density of
the linked experiment relative to it is
\[
 L-n\log2+F(G)-\beta\widehat H_G(x).
\]
Differentiating the log partition function gives a Gaussian Lipschitz
constant at most $C_\beta\sqrt n$. The Gaussian concentration
inequality~\cite[Section~5.4, Theorem~5.6]{BLM2013} and the pressure bound
therefore imply
\[
 F(G)\le n\mathcal P+n\eta
\]
outside an event of probability at most
$2\exp(-c_\beta n\eta^2)$. This uses the ordinary Gaussian marginal
of $G$, which remains unchanged after linking.

Substitute \eqref{eq:3}, \eqref{eq:7}, and
$p(V)^2\le C_\beta e^{-V}$ into the log density. On the events just
proved, it is at most
\[
 C_\beta n\bigl(h+\eta\log(2K)+e^{-V}\bigr)
 =o\bigl(n^{1-l/2}\bigr).
\]
Indeed, $V=l\log n+O(h)$, $\log(2K)=O(\log n)$, and
$0<l\le1/100<1/16$. The union of the exceptional events has
probability at most $\exp(-n^{1-10l})$ for sufficiently large $n$.
Taking the restricted output marginal gives
\eqref{eq:6}. Its good event is defined entirely from the inputs of
this one proposal run.

Finally, the same comparison in expectation proves \eqref{eq:5}
by contraction of relative entropy under marginalization. The
exceptional paths contribute negligibly: $L$ has second moment
$O(n^2(1+V)^2)$, and the supremum of the absolute spin energies has
second moment $O(n^2)$ by Gaussian tails and a union bound. The log
partition function is bounded by $n\log2$ plus that energy supremum.
Cauchy--Schwarz controls all exceptional contributions. The Gram
estimate was already established in the first part of the proof.
\end{proof}
